\section{Phase-velocity slices}\label{sec:slices}

This section diagonalises the operator $L$ of one branch and shows that the diagonalising transform commutes with the nonlinear terms of the kinetic equation. We consider the branch equation (\ref{eq:branch}) with the operator (\ref{eq:L}),
\begin{equation*}
  \partial_tG+\{\Phi,G\}+s\,\gradpar(LG)=0,\qquad
  LG=uG-\frac{uF_M}{\Lambda}\,M_0[G],
\end{equation*}
in which $G(\boldsymbol r,u,t)$ is $G^+$ or $G^-$, with $\Lambda=\Lambda^+$ or $\Lambda^-$ from (\ref{eq:Lambdapm}), or a Laguerre coefficient $G_l$ with $l\ge2$, for which $\Lambda=\infty$. The Alfv\'enic fields $\Phi$ and $\Psi$ are treated as given functions of $\boldsymbol r$ and $t$, and nothing in this section uses the RMHD equations (\ref{eq:rmhd}). Two properties of (\ref{eq:branch}) are used: $L$ acts on $u$ alone, with a coefficient $1/\Lambda$ that is the same at every point and on every scale, and the brackets $\{\Phi,\cdot\}$ and $\{\Psi,\cdot\}$ act in the same way at every $u$.

\subsection{The free-energy product and the factorisation \texorpdfstring{$L=UP$}{L = UP}}\label{sl:product}

Let $\mathcal W$ be the space of real functions $G(u)$ with $\ip{G}{G}<\infty$, where
\begin{equation*}
  \ip{f}{g}=\int\frac{f\,g}{F_M}\,\dd u .
\end{equation*}
Since $M_0[g]=\ip{F_M}{g}$ and $\ip{F_M}{F_M}=1$, the operator $\mathsf Q=F_M\,M_0$ is the orthogonal projector onto $F_M$ in this inner product. The free-energy product is
\begin{equation}\label{eq:fe}
  \fe{f}{g}=\ip{f}{g}-\frac{M_0[f]\,M_0[g]}{\Lambda}=\ip{f}{Pg},\qquad
  P=I-\frac{\mathsf Q}{\Lambda}=I-\frac{F_M\,M_0}{\Lambda} .
\end{equation}
For a branch, $\tfrac12 n_iT_i\fe{G^\pm}{G^\pm}$ averaged over the volume is the free energy $W^\pm$ of the branch in the form of CMM's eq.~(B13). S09 (\S6.2.3) showed that it is conserved, and section~\ref{sec:invariants} returns to this identification. Writing $G=M_0[G]\,F_M+G_\perp$ with $M_0[G_\perp]=0$ gives
\begin{equation}\label{sl:split}
  \fe{G}{G}=\ip{G_\perp}{G_\perp}+p_0\,M_0[G]^2,\qquad p_0=1-\frac{1}{\Lambda} .
\end{equation}
By (\ref{sl:split}) the product is positive definite if and only if $p_0>0$, that is, for $\Lambda\notin[0,1]$, and then
\begin{equation*}
  \min(1,p_0)\,\ip{G}{G}\le\fe{G}{G}\le\max(1,p_0)\,\ip{G}{G},
\end{equation*}
so that $(\mathcal W,\fe{\cdot}{\cdot})$ is a Hilbert space whose norm is equivalent to that of $\ip{\cdot}{\cdot}$. Since $\Lambda^+>1$ and $\Lambda^-<0$ for all $\beta_i$ and $\tau/Z$ (section~\ref{sec:model}), and $p_0=1$ for $\Lambda=\infty$, this covers every compressive sector of KRMHD. At $\Lambda=1$ the product is degenerate, with null vector $F_M$. For $0<\Lambda<1$ it has exactly one negative direction, since $\fe{F_M}{F_M}=p_0<0$ while $\fe{G}{G}>0$ for every non-zero $G$ with $M_0[G]=0$.

Let $U$ be multiplication by $u$, $(UG)(u)=uG(u)$, with domain $D(U)=\{G\in\mathcal W:\ uG\in\mathcal W\}$; it is self-adjoint in $(\mathcal W,\ip{\cdot}{\cdot})$. Since $UPG=uG-uF_M\,M_0[G]/\Lambda$, the branch operator factorises as
\begin{equation}\label{eq:LUP}
  L=UP .
\end{equation}
Because $P$ and $U$ are symmetric in $\ip{\cdot}{\cdot}$, (\ref{eq:fe}) and (\ref{eq:LUP}) give
\begin{equation*}
  \fe{f}{Lg}=\ip{Pf}{UPg}
  =\int\frac{u\,f\,g}{F_M}\,\dd u-\frac{M_0[f]\,M_1[g]+M_1[f]\,M_0[g]}{\Lambda},
\end{equation*}
where the last step uses $M_1[F_M]=0$. The right-hand side is symmetric in $f$ and $g$, so $L$ is symmetric in the free-energy product. For $\Lambda\notin[0,1]$ it is also self-adjoint, with domain $D(U)$: the map $G\mapsto P^{1/2}G$ is unitary from $(\mathcal W,\fe{\cdot}{\cdot})$ onto $(\mathcal W,\ip{\cdot}{\cdot})$ and carries $L$ into $P^{1/2}UP^{1/2}$, which differs from $U$ by a bounded symmetric operator of rank two (appendix~\ref{app:spectral}). Since $P$ is invertible for $\Lambda\neq1$, (\ref{eq:LUP}) also gives
\begin{equation}\label{sl:PLinv}
  PL^{-1}=U^{-1},\qquad\text{that is,}\qquad
  \fe{f}{L^{-1}g}=\int\frac{f\,g}{u\,F_M}\,\dd u\quad\text{for } g/u\in\mathcal W .
\end{equation}
The form (\ref{sl:PLinv}) contains no $\Lambda$; section~\ref{sec:invariants} uses it for CMM's phase-space helicity.

In the Hermite representation of section~\ref{sec:model}, let $e_m=H_m(u)F_M(u)/\sqrt{2^mm!}$, so that $G=\sum_mg_me_m$. The $e_m$ are orthonormal in $\ip{\cdot}{\cdot}$, the matrix of $P$ is the weight matrix $\Pw=\mathrm{diag}(p_m)$, and the matrix of $U$ is the Hermite Jacobi matrix $J$, with zero diagonal and $J_{m,m+1}=J_{m+1,m}=\sqrt{(m+1)/2}$. The product (\ref{eq:fe}) of $f=\sum_mf_me_m$ and $g=\sum_mg_me_m$ becomes $\fe{f}{g}=\sum_mp_mf_mg_m$, and the factorisation (\ref{eq:LUP}) becomes $A=J\Pw$ for the ladder matrix of section~\ref{sec:model}, which is why $\Pw A=\Pw J\Pw$ is symmetric.

\subsection{Resolvent and dielectric function}\label{sl:resolventsec}

For $\mathrm{Im}\,z>0$, the equation $(z-L)X=G$ gives $X=[G-uF_M\,M_0[X]/\Lambda]/(z-u)$, and taking $M_0$ of both sides determines $M_0[X]$:
\begin{equation}\label{sl:resolvent}
  (z-L)^{-1}G=\frac{1}{z-u}\left[G-\frac{uF_M}{\Lambda\,\eps(z)}\int\frac{G(v)}{z-v}\,\dd v\right],
\end{equation}
with the dielectric function
\begin{equation}\label{eq:eps}
  \eps(\zeta)=1-\frac{1}{\Lambda}\int\frac{vF_M(v)}{v-\zeta}\,\dd v=1-\frac{1+\zeta Z(\zeta)}{\Lambda},
  \qquad Z(\zeta)=\int\frac{F_M(x)}{x-\zeta}\,\dd x\quad(\mathrm{Im}\,\zeta>0).
\end{equation}
Here $Z$ is the plasma dispersion function, and $\eps=0$ is the dispersion relation of the branch (section~\ref{sec:model}). We use $\eps$ for the function (\ref{eq:eps}) in the upper half plane and for its analytic continuation into the lower half plane; the zeros of the continuation are the Landau roots. For $\mathrm{Im}\,z<0$ the resolvent is given by (\ref{sl:resolvent}) with $\eps(z)$ replaced by $\overline{\eps(\bar z)}$, which is the value of the integral in (\ref{eq:eps}) at $z$ and differs from the analytic continuation of $\eps$. On the real axis the Plemelj formula gives
\begin{equation}\label{sl:epsRI}
  \eps(u+\ii0)=\eps_R(u)+\ii\,\eps_I(u),\qquad
  \eps_R=1-\frac{1+u\,\mathrm{Re}\,Z(u)}{\Lambda},\qquad
  \eps_I=-\frac{\pi\,u\,F_M(u)}{\Lambda},
\end{equation}
where $\mathrm{Re}\,Z(u)=\PVint F_M(x)\,(x-u)^{-1}\,\dd x$ is minus twice the Dawson function. With $\mathsf H[f](u)=\pi^{-1}\PVint f(v)\,(v-u)^{-1}\,\dd v$, which has the Fourier multiplier $\ii\,\mathrm{sgn}\,k$ and is minus the transform usually called the Hilbert transform (we call $\mathsf H$ the Hilbert transform below), the real and imaginary parts are related by $\eps_R-1=\mathsf H[\eps_I]$.

By (\ref{sl:resolvent}), the resolvent exists at every $z$ in the upper half plane at which $\eps(z)\neq0$, and at the complex conjugates of these points. Appendix~\ref{app:spectral} shows that the non-real eigenvalues of $L$ are the zeros of $\eps$ in the upper half plane and their complex conjugates, and that a real eigenvalue can occur only at $u=0$ and only for $\Lambda=1$. The zeros of $\eps$ in the upper half plane are counted by the argument of \citet{Penrose1960}, which appendix~\ref{app:spectral} carries out for this $\eps$: as $u$ runs over the real line, the curve $1+uZ(u+\ii0)$ winds once around every point of the interval $(0,1)$ and around no real point outside $[0,1]$. Hence $\eps$ has no zeros in the closed upper half plane for $\Lambda\notin[0,1]$, and exactly one, on the imaginary axis, for $0<\Lambda<1$. In the latter case $L$ has the eigenvalues $\pm\ii\gamma$, where $\gamma>0$ solves $1-\sqrt\pi\,\gamma\,e^{\gamma^2}\mathrm{erfc}(\gamma)=\Lambda$, and the slice decomposition below is incomplete (appendix~\ref{app:spectral}). At $\Lambda=1$, $\eps(0+\ii0)=0$ and $F_M$ is an eigenfunction with eigenvalue zero embedded in the continuous spectrum. The branch $\Lambda^+$ approaches this case at high beta. For $\Lambda=1+\delta$ the least-damped zero of the analytic continuation of $\eps$ is $\zeta=-(\ii\delta/\sqrt\pi)\,[1+O(\delta)]$. Since $\Lambda^+-1\simeq\beta_i^{-1}$ for $\beta_i\gg1$, the corresponding damping rate is $|k_\parallel|v_A/\sqrt{\pi\beta_i}$, the Barnes damping rate \citep{Barnes1966} in the form of S09's eq.~(190).

The spectral density of the branch is
\begin{equation}\label{eq:varrho}
  \varrho(u)=\frac{F_M(u)}{|\eps(u+\ii0)|^2} .
\end{equation}
For $\Lambda\notin[0,1]$ it is positive, smooth and even, with
\begin{equation*}
  \int\varrho\,\dd u=\frac{\Lambda}{\Lambda-1}=\frac{1}{p_0},\qquad
  \varrho(0)=\frac{F_M(0)}{p_0^2},\qquad
  \frac{\varrho(u)}{F_M(u)}\to1\quad(|u|\to\infty),
\end{equation*}
and $\varrho=F_M$ for $\Lambda=\infty$. At $\Lambda=1$ it behaves as $\pi^{-3/2}u^{-2}$ near $u=0$ and is not integrable; this is the trace of the embedded eigenvalue. Figure~\ref{fig:spectrum} shows $\varrho$ for $\Lambda=\pm\sqrt5$.

\subsection{Case--Van Kampen modes and the slice transform}\label{sl:transformsec}

For real $u$, the equation $(L-u)F=0$ has the distributional solutions
\begin{equation}\label{sl:Fu}
  F_u(v)=\frac{1}{\Lambda}\,\mathrm P\frac{vF_M(v)}{v-u}+\eps_R(u)\,\delta(v-u),\qquad
  M_0[F_u]=1,\qquad M_1[F_u]=u,
\end{equation}
which are the Case--Van Kampen modes of the branch \citep{vanKampen1955,Case1959}. The normalisation $M_0[F_u]=1$ fixes the weight of the $\delta$-function. Then $(v-u)F_u=vF_M/\Lambda$, and $(L-u)F_u=(v-u)F_u-vF_M\,M_0[F_u]/\Lambda=0$. Each mode consists of a beam of particles at $v=u$, with weight $\eps_R(u)$, together with a principal-value part, which is the response of the Maxwellian, through the Landau coupling, to the density $M_0[F_u]=1$ of the mode moving at the phase velocity $u$. For $\Lambda=\infty$ the principal-value part vanishes and $F_u=\delta(v-u)$.

The slice transform is the projection on these modes in the free-energy product,
\begin{equation}\label{eq:transform}
  \Tr[G](u)=\fe{F_u}{G}
  =\frac{\eps_R(u)\,G(u)}{F_M(u)}+\frac{u}{\Lambda}\,\PVint\frac{G(v)}{v-u}\,\dd v
  =\frac{\eps_R\,G-\eps_I\,\mathsf H[G]}{F_M}(u).
\end{equation}
Since $LF_u=uF_u$ and $L$ is symmetric,
\begin{equation}\label{sl:TL}
  \Tr[LG](u)=u\,\Tr[G](u)\quad\text{for }G\in D(U),\qquad \Tr[F_M](u)=1-\frac{1}{\Lambda}\quad\text{for every }u ;
\end{equation}
appendix~\ref{app:spectral} verifies the first identity directly from (\ref{eq:transform}). The unit-weight slice amplitude is
\begin{equation}\label{eq:aunit}
  a_u=\sqrt{\varrho(u)}\;\Tr[G](u)=\frac{\eps_R\,G-\eps_I\,\mathsf H[G]}{\sqrt{F_M}\;|\eps|}(u),
\end{equation}
with $\varrho$ from (\ref{eq:varrho}), and we write $\mathcal A$ for the linear map $G\mapsto a$. For $\Lambda=\infty$, $a_u=G(u)/\sqrt{F_M(u)}$: the Laguerre coefficients $G_l$ with $l\ge2$ are already slices.

The modes (\ref{sl:Fu}) are those of \citet{vanKampen1955} and \citet{Case1959}, and the reconstruction formula (\ref{sl:inverse}) below is the G-transform $\mathcal G[c]=\eps_Rc+\eps_I\mathsf H[c]$ of \citet{MorrisonPfirsch1992}, \citet{Morrison1994} and \citet{Morrison2000}. \citet{Morrison2000} proved for the linearised Vlasov--Poisson system at fixed wavenumber that this transform is invertible when the Penrose condition holds, and \citet{Heninger2018} applied it to electrostatic drift kinetics (remark~\ref{sl:remlit}). The following theorem gives the corresponding statements for $L$, in the free-energy product and with the weights that make the map unitary.

\begin{theorem}\label{thm:spectrum}
Let $\Lambda\notin[0,1]$, the value $\Lambda=\infty$ included, and let spectral statements refer to the complexification of $\mathcal W$. Then:
\begin{itemize}
\item[(i)] $L$, with domain $D(U)$, is self-adjoint in $(\mathcal W,\fe{\cdot}{\cdot})$;
\item[(ii)] $L$ has no eigenvalues, and its spectrum is the real line, purely absolutely continuous and of multiplicity one;
\item[(iii)] $\mathcal A$ is a unitary map from $(\mathcal W,\fe{\cdot}{\cdot})$ onto $L^2(\mathbb R,\dd u)$, and $\mathcal ALG=u\,\mathcal AG$ for $G\in D(U)$.
\end{itemize}
In particular, for every bounded Borel function $\varphi$,
\begin{equation}\label{sl:parseval}
  \fe{G}{G}=\int a_u^2\,\dd u,\qquad \fe{G}{\varphi(L)G}=\int\varphi(u)\,a_u^2\,\dd u,
\end{equation}
and $G$ is recovered from its slices by
\begin{equation}\label{sl:inverse}
  G=\int\varrho(u)\,\Tr[G](u)\,F_u\,\dd u=\eps_R\,c+\eps_I\,\mathsf H[c],\qquad c=\varrho\,\Tr[G]=\sqrt\varrho\;a .
\end{equation}
\end{theorem}

Appendix~\ref{app:spectral} proves theorem~\ref{thm:spectrum}. In outline, the unitary map $G\mapsto P^{1/2}G$ carries $L$ into $U+(\sqrt{p_0}-1)(\mathsf QU+U\mathsf Q)$, which is $U$ plus a bounded symmetric perturbation of rank two; this gives part~(i). A perturbation of rank two is of trace class, so by the Kato--Rosenblum theorem \citep[][chapter~X]{Kato1966} the absolutely continuous part of $L$ is unitarily equivalent to $U$, whose spectrum is the real line with multiplicity one. Because $\eps$ has no zeros in the closed upper half plane, $1/\eps$ is continuous up to the real axis, and the resolvent form $\fe{G}{(z-L)^{-1}G}$ has continuous boundary values on the real axis for a dense set of $G$. This excludes eigenvalues and singular continuous spectrum, which gives part~(ii). Stone's formula then gives the spectral measure of $L$ in the state $G$ as $\varrho\,\Tr[G]^2\,\dd u=a_u^2\,\dd u$, which is (\ref{sl:parseval}), so that $\mathcal A$ is an isometry, and the construction of a preimage for every smooth $a$ with compact support shows that $\mathcal A$ is onto. The reconstruction formula (\ref{sl:inverse}) uses the absence of zeros of $\eps$ in the closed upper half plane. The formula (\ref{eq:aunit}) holds for every $G\in\mathcal W$, with $\mathsf H$ the Hilbert transform on $L^2(\mathbb R,\dd u)$. For complex $G$, such as a Fourier coefficient, $a_u^2$ in (\ref{sl:parseval}) is replaced by $|a_u|^2$.

Theorem~\ref{thm:spectrum} fixes the weights: $\mathcal G$ is unitary from $L^2(\mathbb R,\dd u/\varrho)$ onto $(\mathcal W,\fe{\cdot}{\cdot})$, with inverse $G\mapsto\varrho\,\Tr[G]$, and the corresponding map onto $L^2(\mathbb R,\dd u)$ is $\mathcal A=\varrho^{-1/2}\mathcal G^{-1}$, that is, $G\mapsto\sqrt\varrho\,\Tr[G]$. In terms of the modes, part~(iii) contains the completeness relation $G=\int\varrho\,\Tr[G]\,F_u\,\dd u$ and the orthogonality relation $\fe{F_u}{F_{u'}}=\delta(u-u')/\varrho(u)$. By (\ref{sl:parseval}), the slices divide the free energy of a branch among phase velocities, with no cross terms between different $u$.

\subsection{The slice theorem}\label{sl:slicesec}

\begin{theorem}\label{thm:slice}
Assume that
\begin{itemize}
\item[(i)] $\Lambda\notin[0,1]$, the value $\Lambda=\infty$ included;
\item[(ii)] $\Phi(\boldsymbol r,t)$ and $\Psi(\boldsymbol r,t)$ are real functions on $\mathbb T^3\times[0,T]$, where $\mathbb T^3$ is a periodic box, and their gradients are continuous and Lipschitz in $\boldsymbol r$, uniformly in $t$; otherwise $\Phi$ and $\Psi$ are arbitrary, and in particular they need not satisfy (\ref{eq:rmhd}): they may obey resistive or forced RMHD, or be prescribed;
\item[(iii)] $G$ and $UG$ are continuously differentiable functions of $(\boldsymbol r,t)$ with values in $\mathcal W$.
\end{itemize}
Then $G$ solves (\ref{eq:branch}) if and only if, at every $(\boldsymbol r,t)$, the slice amplitudes $a_u(\boldsymbol r,t)$ of (\ref{eq:aunit}) satisfy
\begin{equation}\label{eq:slice}
  \partial_ta_u+\{\Phi+su\Psi,\,a_u\}+su\,\partial_za_u=0
\end{equation}
for almost every $u\in\mathbb R$; this is the advection equation
\begin{equation}\label{eq:Uu}
  \partial_ta_u+\boldsymbol U_u\cdot\nabla a_u=0,\qquad
  \boldsymbol U_u=\zhat\times\nabla_\perp(\Phi+su\Psi)+su\,\zhat=\uE+su\,\bhat,\qquad
  \nabla\cdot\boldsymbol U_u=0 .
\end{equation}
The same holds with $a_u$ replaced by $\Tr[G](u)$. By (ii) the velocity $\boldsymbol U_u$ is Lipschitz in $\boldsymbol r$, so that each slice is transported along the characteristics of $\boldsymbol U_u$.
\end{theorem}

\begin{proof}
By theorem~\ref{thm:spectrum}(iii), $\mathcal A$ is a bounded, injective linear map that acts on $u$ only, through a kernel that does not depend on $\boldsymbol r$ or $t$. It therefore commutes with $\partial_t$, $\partial_x$, $\partial_y$ and $\partial_z$ applied to continuously differentiable $\mathcal W$-valued functions, and with multiplication by functions of $(\boldsymbol r,t)$; hence it commutes with $\{\Phi,\cdot\}$, with $\{\Psi,\cdot\}$ and with $\gradpar=\partial_z+\{\Psi,\cdot\}$. By (iii), $LG=UG-uF_M\,M_0[G]/\Lambda$ is continuously differentiable, and $\mathcal ALG=u\,\mathcal AG$ then gives
\begin{equation*}
  \mathcal A\big[\partial_tG+\{\Phi,G\}+s\,\gradpar(LG)\big]
  =\partial_ta+\{\Phi,a\}+s\,u\,\gradpar a,
\end{equation*}
with $a=\mathcal AG$ and $u$ acting by multiplication. Since $u\,\gradpar a_u=u\,\partial_za_u+\{u\Psi,a_u\}$, the right-hand side is the left-hand side of (\ref{eq:slice}). Because $\mathcal A$ is injective, the expression in square brackets vanishes at $(\boldsymbol r,t)$ if and only if the right-hand side vanishes there as an element of $L^2(\mathbb R,\dd u)$, that is, for almost every $u$. The identity $\{\chi,a\}=(\zhat\times\nabla_\perp\chi)\cdot\nabla a$ gives (\ref{eq:Uu}), and $\nabla\cdot(\zhat\times\nabla_\perp\chi)=0$ gives $\nabla\cdot\boldsymbol U_u=0$. Since $\varrho(u)$ is positive and finite, dividing (\ref{eq:slice}) by $\sqrt{\varrho(u)}$ gives the statement for $\Tr[G](u)$.
\end{proof}

When $a_u$ and its derivatives are continuous in $u$, as they are for a finite Hermite expansion of $G$, (\ref{eq:slice}) holds for every $u$. Each slice is a passive scalar carried by the divergence-free velocity $\boldsymbol U_u$, and slices with different $u$ do not interact. The statement is three-dimensional and includes the part $s\{\Psi,LG\}$ of the streaming that follows the perturbed field lines. The velocity $\boldsymbol U_u$ is that of the guiding centre of a free ion that drifts with $\uE$ and streams along the perturbed field line $\bhat=\zhat+\boldsymbol b_\perp$ at the constant parallel velocity $v_\parallel=u\vth$, which is $su$ in units of $v_A$. The parallel electric field and the mirror force, which couple the ions to $\delta n$ and $\delta B$ and enter (\ref{eq:branch}) through the Landau coupling $uF_M\,M_0[G]/\Lambda$, do not appear in (\ref{eq:slice}). The Landau coupling only dresses the slices: its effect is contained in the relation (\ref{eq:aunit}) between $a_u$ and $G$, by which the slice at $u$ is the projection of $G$ on the mode $F_u$, the beam of particles at $v=u$ together with the response (\ref{sl:Fu}) of the Maxwellian to it. The distribution $G(\boldsymbol r,u,t)$ at fixed $u$ is not a passive scalar, since the Landau coupling links it to all other velocities; the dressed combination $a_u$ is. In a straight field the slice streams at $su$, the parallel phase velocity of the mode $F_u$, and we call $u$ the phase velocity of the slice.

The proof uses two properties of the branch equation. The Landau coupling is linear in $G$, local in $\boldsymbol r$, and has the coefficient $1/\Lambda$, which does not depend on the wavenumber, so a single transform diagonalises $L$ at every wavenumber; and the brackets act in the same way at every $u$. Neither property holds for the Vlasov--Poisson system, in which the analogue of $\Lambda$ is $-(k\lambda_D)^2$ (section~\ref{sec:model}) and the nonlinear term acts on velocity. That the slices are also passive is a further property of KRMHD: $\Phi$ and $\Psi$ obey (\ref{eq:rmhd}), which does not involve $G$. When the advecting fields depend on the distribution, as the electrostatic potential does in the drift-kinetic system of \citet{Heninger2018}, the transformed components obey (\ref{eq:slice}) with $\Psi=0$ and with a potential that depends on them, and they remain coupled through it. Passivity is not needed for the energy of each slice to be conserved (section~\ref{sec:invariants}), but it is needed for the slices to evolve independently of one another. Both the decomposition and the passivity are properties of the limit $k_\perp\rho_i\to0$. At finite $k_\perp\rho_i$ the field relations contain gyroaveraging factors, so that the effective $\Lambda$ depends on $k_\perp$ and the advecting potential on $v_\perp$, and the Alfv\'enic fields depend on the compressive fluctuations; section~\ref{sec:invariants} states what can be proved there.

The slice at $u=0$ is advected by $\uE$ alone. If $G$ is H\"older continuous in $u$ near $u=0$, (\ref{eq:transform}) gives $\Tr[G](0)=p_0\,G(0)/F_M(0)$, so this slice is the distribution at $v_\parallel=0$ itself, which obeys $\partial_tG+\{\Phi,G\}=0$ there because $(LG)(0)=0$. This is the KRMHD form of CMM's eq.~(2.22), by which in gyrokinetics the distribution at $v_\parallel=0$ is advected by the gyroaveraged potential alone; section~\ref{sec:invariants} relates the energy of this slice to the imaginary part of CMM's complex helicity.

In terms of the Elsasser potentials $z^\pm=\Phi\pm\Psi$, the stream function of the slice at $u$ is
\begin{equation}\label{eq:elsasserpot}
  \Phi+su\Psi=\tfrac12\big[(1+su)\,z^++(1-su)\,z^-\big] .
\end{equation}
By (\ref{eq:elsasserpot}), the slices with $su=\pm1$, that is, with $v_\parallel=\pm v_A$, are advected by a single Elsasser field (proposition~\ref{prop:elsasser}).

The two branches and the sector $\Lambda=\infty$ have different transforms, because their values of $\Lambda$ differ, but the velocity (\ref{eq:Uu}) depends only on $u$. Applying theorem~\ref{thm:slice} to each of them, with $a^\pm_u$ the slices of $G^\pm$ and $a^{(l)}_u=G_l(u)/\sqrt{F_M(u)}$ for $l\ge2$, gives
\begin{equation}\label{eq:sameflow}
  \big(\partial_t+\boldsymbol U_u\cdot\nabla\big)\big(a^+_u,\,a^-_u,\,a^{(2)}_u,\,a^{(3)}_u,\dots\big)=0 .
\end{equation}
By (\ref{eq:sameflow}), the compressive sector of KRMHD is a continuum of passive scalars labelled by phase velocity, and at each $u$ the slices of all sectors are carried by the same flow. Section~\ref{sec:truncated} shows that the truncated Hermite hierarchy integrated by GANDALF decomposes in the same way, with the slices at the eigenvalues of the truncated ladder matrix. We checked that decomposition to round-off against GANDALF's right-hand side for $\Lambda=\pm\sqrt5$ and $\beta_i=0.3$, $0.5$ and $1$, with both closures at $\beta_i=0.3$ and $1$, in states in which the brackets were up to several times larger than the streaming term. In part of these checks the brackets were evaluated independently of GANDALF, and the same comparison with the sign of $\{\Psi,\cdot\}$ reversed fails by an amount of order one (section~\ref{sec:numerics}).

\subsection{Dissipation, collisions and forcing}\label{sl:forcingsec}

Theorem~\ref{thm:slice} concerns the collisionless and unforced branch equation. Consider instead
\begin{equation*}
  \partial_tG+\{\Phi,G\}+s\,\gradpar(LG)=-\mathcal BG-\mathcal CG+F,
\end{equation*}
where $\mathcal B$ is a linear operator that acts on $\boldsymbol r$ alone and in the same way at every $u$, $\mathcal C$ is a linear operator that acts on $u$ alone and in the same way at every $\boldsymbol r$, and $F(\boldsymbol r,u,t)$ is a source. Under the hypotheses of theorem~\ref{thm:slice}, and with $\mathcal BG$, $\mathcal CG$ and $F$ continuous functions of $(\boldsymbol r,t)$ with values in $\mathcal W$, the argument of its proof gives
\begin{equation}\label{eq:sliceforcing}
  \partial_ta_u+\boldsymbol U_u\cdot\nabla a_u=-\mathcal B\,a_u-\int\hat{\mathcal C}(u,u')\,a_{u'}\,\dd u'+\sqrt{\varrho(u)}\;\Tr[F](u),
  \qquad
  \hat{\mathcal C}(u,u')=\sqrt{\varrho(u)\,\varrho(u')}\;\fe{F_u}{\mathcal CF_{u'}} ,
\end{equation}
where $\hat{\mathcal C}$ is the kernel of $\mathcal A\,\mathcal C\,\mathcal A^{-1}$.

An operator $\mathcal B$ commutes with $\mathcal A$ and appears unchanged in every slice equation. Damping at a rate that depends only on $k_\perp$ or $k_z$, applied with the same coefficient to every Hermite moment, is of this kind; GANDALF applies its hyper-resistivity and its parallel damping in this way \citep{Kanekar2025}. Such terms damp each slice and do not couple slices. We checked (\ref{eq:sliceforcing}) to round-off with a damping term proportional to $k_\perp^4$.

An operator $\mathcal C$ that depends on velocity couples the slices through the kernel $\hat{\mathcal C}$. \citet{Heninger2018}, whose subject is the G-transform of kinetic equations with collisions, state that for electrostatic drift kinetics the transform changes the collision operator only by a small term. The kernel is explicit for an operator that is diagonal in the Hermite basis, $\mathcal Ce_m=\nu_me_m$. Examples are the Lenard--Bernstein operator in $v_\parallel$, $\mathcal CG=-\nu\,\partial_u(uG+\tfrac12\partial_uG)$, for which $\nu_m=\nu m$, and GANDALF's hypercollisions, for which $\nu_m=\nu(m/M)^n$ for $m\ge2$ and $\nu_0=\nu_1=0$. Let $r_m(u)$ be the $m$th Hermite coefficient of the mode $F_u$. It is a polynomial of degree $m$ (section~\ref{sec:truncated}), and by (\ref{sl:Fu})
\begin{equation*}
  r_0(u)=M_0[F_u]=1,\qquad r_1(u)=\sqrt2\,M_1[F_u]=\sqrt2\,u .
\end{equation*}
Since $\fe{f}{g}=\sum_mp_mf_mg_m$,
\begin{equation}\label{sl:Ten}
  \Tr[e_m](u)=p_m\,r_m(u).
\end{equation}
The functions $e_m/\sqrt{p_m}$ form an orthonormal basis of $(\mathcal W,\fe{\cdot}{\cdot})$, so by theorem~\ref{thm:spectrum} their images $\phi_m=\mathcal A[e_m/\sqrt{p_m}]=\sqrt{p_m\varrho}\;r_m$ form an orthonormal basis of $L^2(\mathbb R,\dd u)$, and
\begin{equation*}
  \hat{\mathcal C}(u,u')=\sum_m\nu_m\,\phi_m(u)\,\phi_m(u') .
\end{equation*}
This kernel is $\nu\,\delta(u-u')$ when $\nu_m=\nu$ for every $m$, and it couples different slices otherwise. For GANDALF's hypercollisions with $n=6$, $M=24$ and $\Lambda=\pm\sqrt5$, the largest off-diagonal element of their matrix in the unit-weight slice variables of the truncated hierarchy (section~\ref{sec:truncated}) is $0.40$ times the largest diagonal element, so the coupling is not small. Collisions are the only term in (\ref{eq:sliceforcing}) that exchanges free energy between slices.

Forcing of the Hermite moment $g_m$ by a field $f(\boldsymbol r,t)$ is the source $F=f\,e_m$, and by (\ref{sl:Ten}) it gives the slice at $u$ the source $\sqrt{\varrho(u)}\,p_m\,r_m(u)\,f$. For the two lowest moments,
\begin{equation}\label{sl:sources}
  F=f\,e_0:\quad \Big(1-\frac{1}{\Lambda}\Big)\sqrt{\varrho(u)}\;f,\qquad\qquad
  F=f\,e_1:\quad \sqrt2\,u\,\sqrt{\varrho(u)}\;f .
\end{equation}
Forcing of $g_0=M_0[G]$ therefore puts the same spatial pattern $f$ into every slice, with an amplitude proportional to $\sqrt\varrho$. For forcing that is white in time, has zero mean and is independent of the state, the free energy injected per unit $u$ is proportional to $\varrho(u)$, with the normalised density $p_0\varrho(u)$, which integrates to one: injection at $m=0$ is uniform with respect to the spectral measure $\varrho\,\dd u$. Forcing of $g_1$ gives the slice at $u$ the amplitude $\sqrt2\,u\sqrt{\varrho(u)}$, which is odd in $u$. The slices at $u$ and $-u$ receive the same pattern with opposite signs, and the injected free energy has the density $2u^2\varrho(u)$, which integrates to $\fe{e_1}{e_1}=1$. In general, such forcing of $g_m$ injects free energy with the normalised density $\phi_m^2=p_m\varrho\,r_m^2$. If $g_0$ and $g_1$ are forced by zero-mean fields $f_0$ and $f_1$ that are white in time and independent of the state, with $\langle f_i(\boldsymbol r,t)\,f_j(\boldsymbol r,t')\rangle=c_{ij}\,\delta(t-t')$, the mean rate at which $a_u^2$ is injected is
\begin{equation*}
  \varrho(u)\,\big[p_0^2\,c_{00}+2\sqrt2\,p_0\,c_{01}\,u+2u^2\,c_{11}\big].
\end{equation*}
The cross term is odd in $u$. It injects no free energy, but it injects $\int u\,a_u^2\,\dd u=\fe{G}{LG}$, whose volume average is $2\Gamma^\pm/\vth$ for CMM's invariant $\Gamma^\pm$ (section~\ref{sec:invariants}), at the rate $\sqrt2\,p_0\,c_{01}$ per unit volume (see also section~\ref{sec:symmetry}).

\begin{remark}\label{sl:remlit}
Parts of this structure are known in other settings. For the sector $\Lambda=\infty$, theorem~\ref{thm:slice} restates S09's eqs~(210) and~(212) (section~\ref{sec:model}), at each $v_\parallel$ for $a_u=G_l(u)/\sqrt{F_M(u)}$. What the transform adds is the treatment of the two branches, whose Landau coupling it absorbs into the slices. \citet[section~IV]{Heninger2018} applied the G-transform to the electrostatic drift-kinetic ion equation with Boltzmann electrons in a slab \citep[see also][chapter~13]{Heninger2020}. There quasineutrality also makes the kernel independent of the wavenumber, and their transform is the one used here with $\Lambda=-T_i/(ZT_e)$, which is the limit of $\Lambda^-$ as $\beta_i\to0$. In their system the transformed components stream along the uniform guide field and are carried by one $\boldsymbol E\times\boldsymbol B$ flow, but the potential is proportional to $\int g\,\dd v_\parallel$, so the components remain coupled through it, and there is no streaming along perturbed field lines. The transform has also been used for linear problems at a single wavenumber, among them a forced Landau problem \citep{Plunk2013} and, in the appendix of \citet{PlunkParker2014}, a kinetic-Alfv\'en-wave model with a fluctuating parallel vector potential. \citet{Plunk2015} used a transform that commutes with $\boldsymbol E\times\boldsymbol B$ advection to obtain a quadratic invariant at each velocity, and from it the nonlinear stability of a quasi-two-dimensional drift-kinetic model without parallel streaming. \citet{Tassi2015} constructed normal fields of truncated Hermite hierarchies of electromagnetic drift-fluid models, and \citet{Tassi2022} identified their eigenvalues $\lambda_i$ as the zeros of a Hermite polynomial and the fields as proportional to the truncated distribution function evaluated at $v_\parallel=\lambda_i$. In two dimensions each field is a Lagrangian invariant of the flow with stream function $\phi-\lambda_i(v_{\mathrm{th}e}/c)A_\parallel$, the analogue of $\Phi+su\Psi$, but $\phi$ and $A_\parallel$ depend on all the fields. In three dimensions the parallel gradients of these potentials also enter the equation of each field, so the fields remain coupled, except, when an odd number of moments is retained, the field with $\lambda_i=0$, which is advected by the $\boldsymbol E\times\boldsymbol B$ flow alone \citep{Tassi2022}. Other two-dimensional precedents are listed in section~\ref{sec:intro}.
\end{remark}
